\begin{definitionbox}{Definition (Difference of Intervals)}
\begin{definition}[Difference of Intervals]
\label{def:difference-of-intervals}
Let $I,J\subseteq\mathbb{R}$ be intervals.  The \emph{difference} of $I$ and
$J$ is the set
\[
  I\setminus J
  :=
  \{x\in\mathbb{R}:x\in I\text{ and }x\notin J\}.
\]
\end{definition}
\end{definitionbox}

\begin{remark*}[Standard quantified statement]
\[
I\setminus J=\{x\in\mathbb{R}:x\in I\land x\notin J\}=I\cap J^c.
\]
\end{remark*}

\begin{remark*}[Predicate reading]
For any $x\in\mathbb{R}$,
\[
x\in I\cup J \Longleftrightarrow (x\in I)\lor(x\in J),
\]
\[
x\in I\cap J \Longleftrightarrow (x\in I)\land(x\in J),
\]
\[
x\in I^c \Longleftrightarrow x\notin I,
\qquad
x\in I\setminus J \Longleftrightarrow (x\in I)\land(x\notin J).
\]
These are set-theoretic operations; no arithmetic operation is being performed
on the endpoints.
\end{remark*}

\begin{remark*}[Interpretation]
The intersection of two intervals is either empty or another interval. For
example, if $a\leq b$ and $c\leq d$, then
\[
[a,b]\cap[c,d]
=
\begin{cases}
[\max\{a,c\},\min\{b,d\}], & \max\{a,c\}\leq \min\{b,d\},\\
\varnothing, & \max\{a,c\}>\min\{b,d\}.
\end{cases}
\]
Endpoint inclusion follows the endpoint rules of the original intervals.


The union of two intervals need not be an interval:
\[
(0,1)\cup(2,3)
\]
has a gap. If two intervals overlap or touch without a missing endpoint, their
union is again an interval, for instance
\[
(0,2]\cup[2,5)=(0,5).
\]
The missing or included endpoints determine whether the joined set is truly
one interval.


The complement of a bounded interval usually splits into two rays:
\[
(a,b)^c=(-\infty,a]\cup[b,\infty),
\qquad
[a,b]^c=(-\infty,a)\cup(b,\infty).
\]
Complements reverse endpoint inclusion: an endpoint excluded from the interval
is included in the complement, and an endpoint included in the interval is
excluded from the complement.


A difference is an intersection with a complement, so it can also split:
\[
(0,5)\setminus[2,3]=(0,2)\cup(3,5).
\]
If the removed interval only cuts off one side, the result may remain a single
interval:
\[
(0,5)\setminus(3,\infty)=(0,3].
\]
The bracket at $3$ appears because $3\notin(3,\infty)$, so it is not removed.


Open sets in $\mathbb{R}$ are built from unions of open intervals. Closed sets
are understood through complements of open sets. Boundaries arise when a point
is not cleanly inside a set or its complement. Thus interval set operations are
the grammar behind the metric topology of the real line.
\end{remark*}

\begin{dependencies}
\begin{itemize}
  \item \hyperref[def:open-interval]{Open Interval}
  \item \hyperref[def:closed-interval]{Closed Interval}
  \item \hyperref[def:left-half-open-interval]{Left-Half-Open Interval}
  \item \hyperref[def:right-half-open-interval]{Right-Half-Open Interval}
  \item \hyperref[def:union-of-intervals]{Union of Intervals}
  \item \hyperref[def:intersection-of-intervals]{Intersection of Intervals}
  \item \hyperref[def:complement-of-interval]{Complement of an Interval}
\end{itemize}
\end{dependencies}
